% MATH 147 course notes.
% Writing standard: handoff.md, section "Presentation standard".
\documentclass[11pt]{article}
% Font: Times + matching math, the qhnotes default.
\usepackage[color]{qhnotes}

% ----- Proof-idea environments (MATH 147 only) -----
% idea:    1--3 sentence strategy, sits directly on top of the proof.
% scratch: backward "how do we find n_eps / choose epsilon" work; not part
%          of the proof, so it gets a grey paper look instead of the proof blue.
\theoremstyle{noteblock}
\newtheorem*{idea}{Idea}
\newtheorem*{scratch}{Scratch work}
\surroundwithmdframed[
  topline=false, bottomline=false, rightline=false, leftline=true,
  linecolor=proofAccent!45, linewidth=1.6pt, backgroundcolor=proofAccent!7,
  skipabove=6pt, skipbelow=0pt,
  innertopmargin=3pt, innerbottommargin=4pt,
  innerleftmargin=8pt, innerrightmargin=4pt
]{idea}
\surroundwithmdframed[
  topline=false, bottomline=false, rightline=false, leftline=true,
  linecolor=black!30, linewidth=1.6pt, backgroundcolor=black!4,
  skipabove=6pt, skipbelow=0pt,
  innertopmargin=3pt, innerbottommargin=4pt,
  innerleftmargin=8pt, innerrightmargin=4pt
]{scratch}

% Document-local shortcuts
\newcommand{\paren}[1]{\left(#1\right)}
\newcommand{\abs}[1]{\left|#1\right|}
\newcommand{\norm}[1]{\left\lVert#1\right\rVert}
\newcommand{\bracks}[1]{\left[#1\right]}
\newcommand{\set}[1]{\left\{#1\right\}}
\newcommand{\R}{\mathbb{R}}
\newcommand{\N}{\mathbb{N}}
\newcommand{\Z}{\mathbb{Z}}
\newcommand{\Q}{\mathbb{Q}}
\newcommand{\C}{\mathbb{C}}
\newcommand{\ds}{\displaystyle}
\newcommand{\blue}[1]{\textcolor{blue}{#1}}
\newcommand{\red}[1]{\textcolor{red}{#1}}
\newcommand{\mypic}[3]{%
  \begin{figure}[htbp]
    \centering
    \includegraphics[width=#3\textwidth]{#1}
    \caption{#2}
  \end{figure}
}

\title{\textbf{Intro to Analysis}\\[0.5em]\large MATH 147}
\author{Qinghao Hu}
\date{\today}
\course{Intro to Analysis}

% Only compile some chapters while drafting:
% \includeonly{chapters/c3}

\begin{document}

\maketitle
\clearpage
\tableofcontents

\section{Course Information}
%!TEX root = ../../note

\subsection{Logistics}

\subsubsection{Contact}
\begin{description}
    \item[Email:] \texttt{mmolino@uwaterloo.ca}
    \item[Website:] \url{www.math.uwaterloo.ca/~mmolino}
    \item[Office hours:] Monday, 3:30--4:30~p.m., and Friday, 10:00--11:00~a.m., in MC 4006.
\end{description}

\subsubsection{Course Structure}
\begin{description}
    \item[Lectures:] Section 002 meets on Monday, Wednesday, and Friday from 12:30--1:20~p.m. in MC 4021.
\end{description}



%!TEX root = ../../note

\section{Foundations of Calculus}
%!TEX root = ../../../note

\subsection{Numbers and Logic}

\subsubsection{What Is Calculus?}
Differential and integral calculus are two branches of calculus.

\begin{description}
    \item[Differential calculus:] The study of slopes and tangent lines.
    \item[Integral calculus:] The study of areas and accumulation.
\end{description}

\subsubsection{From $\N$ to $\Q$}
The number systems discussed in this lecture include
\begin{align*}
    \N &= \set{1,2,3,\ldots}, \\
    \Z &= \set{\ldots,-1,0,1,\ldots}, \\
    \Q &= \set{\frac{m}{n} : m,n\in\Z,\ n\ne 0}.
\end{align*}

The natural numbers are closed under addition. The integers are closed under addition, subtraction, and multiplication, but they are not closed under division; for example, $1/2\notin\Z$. Therefore, we introduce the rational numbers to include quotients of integers with nonzero denominators.

\begin{theorem}[Irrationality of $\sqrt{2}$]
    The number $\sqrt{2}$ is irrational; equivalently, $\sqrt{2}\notin\Q$.
    \insymbols{\neg(\exists m\in\Z)(\exists n\in\N)\,\Bigl[\sqrt2=\frac mn\Bigr]}
\end{theorem}

\begin{idea}
    Write $\sqrt2=m/n$ in lowest terms. Then $m^2=2n^2$ forces $m$ even, and feeding that back forces $n$ even --- so $2$ divides both, contradicting lowest terms.
\end{idea}

\begin{proof}
    Suppose, for contradiction, that $\sqrt{2}\in\Q$. Then we can write
    \begin{equation*}
        \sqrt{2}=\frac{m}{n},
        \qquad m,n\in\Z,
        \qquad n>0,
        \qquad \gcd(m,n)=1.
    \end{equation*}
    Squaring both sides gives
    \begin{equation*}
        2=\frac{m^2}{n^2}
        \qquad\Longrightarrow\qquad
        m^2=2n^2.
    \end{equation*}

    Thus $m^2$ is even. An odd integer has an odd square, so $m$ must be even. Write $m=2\ell$ for some $\ell\in\Z$. Substituting this into the equation above gives
    \begin{align*}
        2n^2 &= 4\ell^2, \\
        n^2 &= 2\ell^2.
    \end{align*}
    Therefore $n^2$ is even, and the same parity argument shows that $n$ is even.

    This means that both $m$ and $n$ are even, contradicting $\gcd(m,n)=1$. Therefore, $\sqrt{2}\notin\Q$.
\end{proof}

\[\boxed{\text{Assuming lowest terms turns the equation into a parity contradiction.}}\]

The irrational numbers are the real numbers that are not rational:
\begin{equation*}
    \R\setminus\Q
    =\set{x\in\R : x\notin\Q}.
\end{equation*}

\begin{example}[Decimal representation]
Every nonnegative real number has a decimal representation of the form
\begin{equation*}
    x=a+\sum_{i=1}^{\infty}\frac{x_i}{10^i},
    \qquad a\in\N\cup\set{0},
    \qquad x_i\in\set{0,1,2,\ldots,9}.
\end{equation*}
Negative real numbers are the negatives of such expansions. Some numbers have two decimal representations, for example $1.999\cdots=2.000\cdots$.
\end{example}

\begin{example}[Recurring decimal]
    Is the statement $1.999\cdots = 2$ correct?

    Yes. Let $x=1.999\cdots$. Then
    \begin{align*}
        x &= 1.999\cdots, \\
        10x &= 19.999\cdots, \\
        9x &= 18, \\
        x &= 2.
    \end{align*}
\end{example}

\subsubsection{Statements and Quantifiers}
\begin{definition}[Statement]
    A statement is a sentence that has a definite truth value: it is either true or false.
\end{definition}

The \textbf{universal quantifier} $\forall$ means ``for all.'' For example,
\begin{example}
    The statement
    \begin{equation*}
        (\forall x \in \R)\,[x^2-x\geq 0]
    \end{equation*}
    is false; for instance, it fails when $x=\frac12$.
\end{example}

The \textbf{existential quantifier} $\exists$ means ``there exists.'' For example,
\begin{example}
    The statement
    \begin{equation*}
        (\exists x \in \R)\,[x^2-x\geq 0]
    \end{equation*}
    is true; for instance, $x=2$ satisfies the inequality.
\end{example}

\begin{notation}
    Each quantifier is written in parentheses, and the statement it governs in square brackets, read left to right:
    \[
        (\forall x\in\R)(\exists y\in\R)\,[x^3-y^3=1]
    \]
    reads ``for every $x\in\R$ there exists $y\in\R$ such that $x^3-y^3=1$.'' Negation moves inward and swaps each quantifier:
    \[
        \neg(\forall x)(\exists y)\,[P(x,y)] \iff (\exists x)(\forall y)\,[\neg P(x,y)].
    \]
    From here on, every theorem-like statement ends with its form \emph{In symbols}: quantifiers and connectives only, using notation defined before that statement.
\end{notation}

\begin{proposition}
    For every $x\in\R$, there exists $y\in\R$ such that
    \begin{equation*}
        x^3-y^3=1.
    \end{equation*}
    \insymbols{(\forall x\in\R)(\exists y\in\R)\,[x^3-y^3=1]}
\end{proposition}

\begin{proof}
    Let $x\in\R$ and take $y=\sqrt[3]{x^3-1}$. Then
    \begin{equation*}
        x^3-y^3=x^3-(x^3-1)=1.
    \end{equation*}
\end{proof}

Here the choice of $y$ is allowed to depend on $x$: for each $x$ we produce a (possibly different) $y$ that makes the equation hold. This is exactly what the quantifier order $(\forall x)(\exists y)$ permits.

\begin{proposition}
    The statement
    \begin{equation*}
        (\exists y\in\R)(\forall x\in\R)\,[x^3-y^3=1]
    \end{equation*}
    is false.
    \insymbols{\neg(\exists y\in\R)(\forall x\in\R)\,[x^3-y^3=1]}

    Its negation is
    \begin{equation*}
        (\forall y\in\R)(\exists x\in\R)\,[x^3-y^3\neq 1].
    \end{equation*}
    To prove the negation, let $y\in\R$ and take $x=y$. Then
    $x^3-y^3=0\neq 1$.
\end{proposition}

\begin{remark}
    In this second statement, the quantifier order is reversed: $(\exists y)(\forall x)$ demands a \emph{single} $y$ that works simultaneously for every $x\in\R$. This is a much stronger requirement than $(\forall x)(\exists y)$, which only asks that a (possibly different) $y$ exist for each individual $x$. The first proposition shows such $y$'s exist, but they genuinely depend on $x$ (namely $y=\sqrt[3]{x^3-1}$), so no single choice of $y$ can work for all $x$ at once. This is why swapping the order of $\forall$ and $\exists$ can change a true statement into a false one, and it is a caution to always read quantified statements in the order they are written.
\end{remark}

%!TEX root=../../note

\subsection{Mathematical Induction}
\begin{definition}[Principle of Mathematical Induction]
    Let $P(n)$ be a statement depending on $n\in\N$. If the following two statements are true:
    \begin{enumerate}[label=(\arabic*)]
        \item $P(1)$ is true;
        \item for every $k\in\N$, if $P(k)$ is true, then $P(k+1)$ is true;
    \end{enumerate}
    then
    \begin{equation*}
        \forall n\in\N,\quad P(n)
    \end{equation*}
    is true.
\end{definition}

\begin{remark}[How to use induction in practice]
    To prove that $P(n)$ holds for every $n\in\N$ by induction, it suffices to carry out two steps: first, verify the \emph{base case} $P(1)$ directly; second, carry out the \emph{inductive step} by fixing an arbitrary $k\in\N$, assuming $P(k)$ holds (the \emph{inductive hypothesis}), and using this assumption to deduce $P(k+1)$. Once both steps are complete, the principle of mathematical induction guarantees $P(n)$ for all $n\in\N$ — no further verification of individual cases is needed.
\end{remark}

Here is an example: 
\begin{proposition}[Sum of the First $n$ Natural Numbers]
    For every $n\in\N$,
    \begin{equation*}
        1+2+\cdots+n=\frac{n(n+1)}{2}.
    \end{equation*}
\end{proposition}


\begin{proof}
    Define $P(n)$ to be the statement
    \begin{equation*}
        1+2+\cdots+n=\frac{n(n+1)}{2}.
    \end{equation*}
    The base case is $n=1$:
    \begin{equation*}
        1=\frac{1(1+1)}{2}.
    \end{equation*}

    For the inductive step, let $k\in\N$ and assume that $P(k)$ is true. Then
    \begin{align*}
        1+2+\cdots+(k+1)
            &= \frac{k(k+1)}{2}+(k+1) \\
            &= \frac{k(k+1)+2(k+1)}{2} \\
            &= \frac{(k+1)(k+2)}{2}.
    \end{align*}
    Thus, $P(k+1)$ is true. By the principle of mathematical induction, the proposition follows.
\end{proof}

The principle of mathematical induction has the following equivalent set formulation.
\begin{proposition}
    Let $S\subseteq\N$. If
    \begin{enumerate}
        \item $1\in S$;
        \item for every $k\in S$, $k+1\in S$;
    \end{enumerate}
    then $S=\N$.
\end{proposition}

\begin{proof}
    Define $P(n)$ to be the statement $n\in S$. The two assumptions give the base case $P(1)$ and the inductive implication $P(k)\Rightarrow P(k+1)$. Hence, by mathematical induction, every $n\in\N$ belongs to $S$. Since $S\subseteq\N$, it follows that $S=\N$.
\end{proof}

\subsubsection{Well-Ordering Principle}
\begin{theorem}[Well-Ordering Principle]
    Every nonempty subset $S$ of $\N$ has a smallest element.
\end{theorem}

\begin{remark}[Proof idea]
    We argue by contradiction: suppose $S$ has no smallest element. We want to show this forces $S=\emptyset$, contradicting the hypothesis that $S$ is nonempty. The natural tool is induction, so we build a set $T$ designed to capture, for each $n$, whether $S$ has been "avoided" up through $n$; showing $T=\N$ by induction will then show $S$ is empty.
\end{remark}

\begin{proof}
    Let $S\subseteq\N$ be nonempty. Suppose, for contradiction, that $S$ has no smallest element. Define
    \begin{equation*}
        T=\set{n\in\N : \set{1,2,\ldots,n}\cap S=\emptyset}.
    \end{equation*}
    That is, $T$ consists of exactly those $n\in\N$ for which none of $1,2,\ldots,n$ lies in $S$ — informally, $n\in T$ means $S$ has "not yet been seen" by the time we reach $n$. If we can show $T=\N$, then in particular no natural number lies in $S$ (since every $n\in\N$ would satisfy $n\in T$, forcing $n\notin S$), giving $S=\emptyset$ and the desired contradiction. We show $T=\N$ by induction on $n$.

    \textbf{Base case.} Since $S$ has no smallest element and $1$ is the smallest element of $\N$, we must have $1\notin S$ (otherwise $1$ would be the smallest element of $S$). Hence $\set{1}\cap S=\emptyset$, so $1\in T$.

    \textbf{Inductive step.} Suppose $k\in T$, i.e., $\set{1,2,\ldots,k}\cap S=\emptyset$. We show $k+1\in T$. Suppose, for contradiction, that $k+1\notin T$; then $\set{1,2,\ldots,k+1}\cap S\neq\emptyset$. Since $\set{1,\ldots,k}\cap S=\emptyset$ by the inductive hypothesis, the only way for $\set{1,\ldots,k+1}\cap S$ to be nonempty is if $k+1\in S$. But then $k+1\in S$ while none of $1,2,\ldots,k$ lie in $S$, so $k+1$ would be the smallest element of $S$ — contradicting our assumption that $S$ has no smallest element. Hence $k+1\in T$.

    By mathematical induction, $T=\N$. This implies that no natural number belongs to $S$, so $S=\emptyset$, contradicting the assumption that $S$ is nonempty. Therefore, $S$ has a smallest element.
\end{proof}

\begin{remark}[Practical use of well-ordering]
    In practice, the well-ordering principle is most often invoked as follows: to show a property $P(n)$ holds for every $n\in\N$, argue by contradiction. If $P(n)$ failed for some $n$, the set of counterexamples $\set{n\in\N : \neg P(n)}$ would be a nonempty subset of $\N$, and the well-ordering principle would guarantee it has a least element $k$. One then derives a contradiction using the minimality of $k$ — typically by showing that $P(k-1)$ must hold (since $k-1$ is smaller than $k$ and hence not a counterexample) and that $P(k-1)$ forces $P(k)$, contradicting $k$ being a counterexample. This is exactly the strategy used in the proposition below.
\end{remark}

\begin{proposition}
    Using the well-ordering principle, prove that for every $n\in\N$,
    \begin{equation*}
        1^2 + 2^2 + \cdots + n^2 = \frac{n\paren{n + 1}\paren{2n + 1}}{6}.
    \end{equation*}
\end{proposition}

\begin{remark}[Proof idea]
    We prove the claim by minimal counterexample. Suppose the formula failed for some $n$; collect all such failures into a set $S=\set{n\in\N:\neg P(n)}$ — this is the natural set to consider because if we can rule out $S$ being nonempty, then $P(n)$ holds for every $n$. Since $S\subseteq\N$, the well-ordering principle (rather than induction directly) is the right tool to extract a concrete least counterexample $k=\min S$. We then show $P(k-1)\Rightarrow P(k)$ by direct computation, so $k$ cannot actually be a counterexample — a contradiction.
\end{remark}

\begin{remark}[Why $\min$, not $\sup$]
    It is worth pausing on why we take the \emph{minimum} of $S$ here rather than some supremum-type argument. The well-ordering principle applies specifically to nonempty subsets of $\N$ (not to general subsets of $\R$), and it guarantees the existence of a smallest element that actually \emph{belongs to} $S$ — this is what allows us to write $k=\min S\in S$ and conclude $\neg P(k)$. This is a stronger statement than merely having an infimum: for a general subset of $\R$, an infimum need not belong to the set, so we could not conclude $\neg P(\inf S)$. Here, because $S\subseteq\N$ and $S$ is nonempty, well-ordering hands us an actual least element $k$ with $\neg P(k)$ true, which is exactly what is needed to test $P$ at $k$ and at $k-1$.

    It is also worth noting that induction and the well-ordering principle are logically equivalent ways of expressing the same underlying fact about $\N$: the proof of the well-ordering principle above used induction, and here we use well-ordering to reprove a statement ($P(n)$ for all $n$) that could equally well have been proved by direct induction on $n$. The two techniques are simply different packagings of the same idea.
\end{remark}

\mdfsetup{nobreak=true}
\begin{proof}
    Let $P(n)$ denote the statement
    \begin{equation*}
        1^2 + 2^2 + \cdots + n^2 = \frac{n(n + 1)(2n + 1)}{6}.
    \end{equation*}
    Suppose, for contradiction, that $P(n)$ fails for some $n$; that is,
    \begin{equation*}
        \neg\paren{\forall n \in \N,\ P(n)}
        \quad\text{equivalently}\quad
        \exists n \in \N,\ \neg P(n).
    \end{equation*}
    Then the set
    \begin{equation*}
        S := \set{n \in \N : \neg P(n)}
    \end{equation*}
    is nonempty. By the well-ordering principle, $S$ has a smallest element $k$. Since $P(1)$ holds, $1\notin S$, so $k > 1$ and hence $k - 1\in\N$.

    By minimality of $k$, we have $k - 1\notin S$, so $P(k - 1)$ is true:
    \begin{equation*}
        1^2 + 2^2 + \cdots + (k - 1)^2
            = \frac{(k - 1)k\paren{2(k - 1) + 1}}{6}
            = \frac{(k - 1)k(2k - 1)}{6}.
    \end{equation*}
    Adding $k^2$ to both sides gives
    \begin{align*}
        1^2 + \cdots + (k - 1)^2 + k^2
            &= \frac{(k - 1)k(2k - 1)}{6} + k^2 \\
            &= \frac{(k - 1)k(2k - 1) + 6k^2}{6} \\
            &= \frac{k\bigl((k-1)(2k-1)+6k\bigr)}{6} \\
            &= \frac{k(2k^2+3k+1)}{6}
             = \frac{k(k+1)(2k+1)}{6},
    \end{align*}
    so $P(k)$ is in fact true. This contradicts $k\in S$. Therefore $S=\emptyset$, and $P(n)$ holds for every $n\in\N$.
\end{proof}
\mdfsetup{nobreak=false}

\begin{remark}[General pattern: minimal counterexample via well-ordering]
    To show $P(n)$ holds for all $n\in\N$: assume the counterexample set $S=\set{n:\neg P(n)}$ is nonempty, let $k=\min S$ (which exists by well-ordering), show $P(k-1)$ holds by minimality of $k$, then show $P(k-1)\Rightarrow P(k)$ — contradicting $k\in S$.
\end{remark}

%!TEX root=../../note
\subsection{The Algebraic and Order Properties of $\R$}
The real numbers $\R$ form a \emph{complete ordered field}. We build up this
description in three stages: the field (algebraic) axioms, the order axioms, and
the completeness axiom.

\begin{definition}[Field]
    A \textbf{field} is a set $F$ equipped with two binary operations
    \[
        + : F \times F \to F,
        \qquad
        \cdot : F \times F \to F,
    \]
    satisfying the following axioms for all $a,b,c \in F$.

    \begin{enumerate}[label=\textbf{F\arabic*.}, leftmargin=2.5em]
        \item \textbf{Associativity of addition.}
        \[
            a + (b + c) = (a + b) + c.
        \]

        \item \textbf{Commutativity of addition.}
        \[
            a + b = b + a.
        \]

        \item \textbf{Additive identity.}
        There exists an element $0 \in F$ such that
        \[
            a + 0 = a.
        \]

        \item \textbf{Additive inverses.}
        For every $a \in F$, there exists an element $-a \in F$ such that
        \[
            a + (-a) = 0.
        \]

        \item \textbf{Associativity of multiplication.}
        \[
            a \cdot (b \cdot c) = (a \cdot b) \cdot c.
        \]

        \item \textbf{Commutativity of multiplication.}
        \[
            a \cdot b = b \cdot a.
        \]

        \item \textbf{Multiplicative identity.}
        There exists an element $1 \in F$, with $1 \neq 0$, such that
        \[
            a \cdot 1 = a.
        \]

        \item \textbf{Multiplicative inverses.}
        For every $a \in F$ with $a \neq 0$, there exists an element $a^{-1} \in F$ such that
        \[
            a \cdot a^{-1} = 1.
        \]

        \item \textbf{Distributive law.}
        \[
            a \cdot (b + c) = a \cdot b + a \cdot c.
        \]
    \end{enumerate}
\end{definition}

\begin{remark}
    The rational numbers $\Q$ also form a field.
\end{remark}

\begin{theorem}
    \begin{enumerate}[label = \alph*)]
        \item If $z, a \in \R$ with $z + a = a$, then $z = 0$.
        \item If $u \in \R$ and $b \neq 0$ with $u \cdot b = b$, then $u = 1$.
        \item If $a \in \R$, then $a \cdot 0 = 0$.
    \end{enumerate}
\end{theorem}

\begin{proof}
    \begin{enumerate}[label = \alph*)]
        \item
        \[
            z = z + 0 = z + (a - a) = (z + a) + (-a) = a + (-a) = 0.
        \]
        \item Multiplying $u \cdot b = b$ by $b^{-1}$,
        \[
            u = u \cdot 1 = u \cdot \paren{b \cdot b^{-1}} = (u \cdot b) \cdot b^{-1} = b \cdot b^{-1} = 1.
        \]
        \item
        \[
            a\cdot0=a\cdot(0+0)=a\cdot0+a\cdot0.
        \]
        Adding $-(a\cdot0)$ to both sides, and using associativity, gives
        $0=a\cdot0$.
    \end{enumerate}
\end{proof}

\begin{definition}[The Order Properties of $\R$]
    There is a nonempty subset $P$ of $\R$, called the set of \textbf{positive real
    numbers}, that satisfies the following properties.
    \begin{enumerate}[label=\textbf{O\arabic*.}, leftmargin=3em]
        \item \textbf{Trichotomy law.}
        For every $a \in \R$, exactly one of the following holds:
        \begin{equation*}
            a = 0, \qquad a \in P, \qquad -a \in P.
        \end{equation*}

        \item \textbf{Closure under addition.}
        If $a, b \in P$, then $a + b \in P$.

        \item \textbf{Closure under multiplication.}
        If $a, b \in P$, then $a \cdot b \in P$.
    \end{enumerate}
\end{definition}

\begin{definition}[Order relation]
    Let $a, b \in \R$.
    \begin{enumerate}
        \item If $a - b \in P$, then we write $a > b$ or $b < a$.
        \item If $a - b \in P \cup \set{0}$, then we write $a \geq b$ or $b \leq a$.
    \end{enumerate}
\end{definition}

\begin{remark}
    $\Q$ is also ordered, so we are still missing something: completeness.
\end{remark}

\begin{definition}[Completeness Axiom]
    If $A \subseteq \R$ is nonempty and bounded above, then $A$ has a least upper
    bound in $\R$.
\end{definition}

\begin{example}
    Prove that if $a \in \R$ and $a \neq 0$, then $a^2 > 0$.
\end{example}

\begin{proof}
    Since $a \neq 0$, the trichotomy law gives two cases.

    \textbf{Case 1: $a \in P$.} By closure under multiplication,
    \begin{equation*}
        a^2 = a \cdot a \in P, \qquad\text{so } a^2 > 0.
    \end{equation*}

    \textbf{Case 2: $a \notin P$.} Since $a \neq 0$, trichotomy gives $-a \in P$, so by closure under multiplication $(-a)(-a) \in P$. It remains to show that $a^2 = (-a)(-a)$:
    \begin{align*}
        a + (-a) = 0
            &\implies (-a)\paren{a + (-a)} = (-a) \cdot 0 = 0 \\
            &\implies (-a)(a) + (-a)(-a) = 0 \\
            &\implies (-a)(-a) = -\paren{(-a)a} = -(-a^2) = a^2.
    \end{align*}
    Hence $a^2 = (-a)(-a) \in P$, so $a^2 > 0$. In either case $a^2 > 0$.
\end{proof}

\subsubsection{The Absolute Value}
\begin{definition}
    Given $a \in \R$, we define the \textbf{absolute value} $\abs{a}$ by
    \[
        \abs{a} =
        \begin{cases}
            a, & a \geq 0, \\
            -a, & a < 0.
        \end{cases}
    \]
\end{definition}
\begin{theorem}
    \begin{enumerate} [label = (\alph*)]
        \item $\abs{ab} = \abs{a}\abs{b}$ for all $a, b \in \R$.
        \item $\abs{a}^2 = a^2$ for all $a \in \R$.
        \item For $a \in \R$ and $c>0$,
        \[\abs{a} \leq c \iff -c \leq a \leq c.\]
        \item For every $a \in \R$,
        \[-\abs{a} \leq a \leq \abs{a}.\]

    \end{enumerate}

\end{theorem}

\begin{proof}
    The definition immediately gives $\abs{-a}=\abs{a}$: if $a\geq0$,
    then $-a\leq0$, and if $a<0$, then $-a>0$. It also gives
    \begin{equation*}
        -\abs{a}\leq a\leq\abs{a},
    \end{equation*}
    proving (d). In the two cases $a\geq0$ and $a<0$, respectively,
    $\abs{a}^2=a^2$ is immediate; this proves (b).

    For (a), there are four sign combinations. If $a,b\geq0$, then
    $ab\geq0$, so $\abs{ab}=ab=\abs a\abs b$. If $a\geq0>b$, then
    $ab\leq0$, so $\abs{ab}=-ab=a(-b)=\abs a\abs b$. The case
    $b\geq0>a$ is the same with the roles reversed. Finally, if $a,b<0$,
    then $ab>0$, and
    \begin{equation*}
        \abs{ab}=ab=(-a)(-b)=\abs a\abs b.
    \end{equation*}

    Finally let $c>0$. Suppose first that $\abs a\leq c$. If $a\geq0$,
    then $\abs a=a$, so $a\leq c$; also $-c<0\leq a$. If $a<0$, then
    $\abs a=-a$, so $-a\leq c$, equivalently $-c\leq a$; also $a<0<c$.
    Hence $-c\leq a\leq c$ in either case.

    Conversely, suppose $-c\leq a\leq c$. If $a\geq0$, then
    $\abs a=a\leq c$. If $a<0$, multiplying $-c\leq a$ by $-1$ reverses
    the inequality and gives $-a\leq c$, hence $\abs a\leq c$. This proves (c).
\end{proof}

\begin{proposition}
    [Triangle Inequality]
    If $a, b \in \R$, then $\abs{a + b } \leq \abs{a } + \abs{b }$.
\end{proposition}

\begin{proof}
    Let $a, b \in \R$. We know that
    \begin{align*}
        -\abs{a } \leq a \leq \abs{a}. \\
        -\abs{b } \leq b \leq \abs{b}.
    \end{align*}

    Thus,
    \begin{align*}
        &-\abs{a } - \abs{b } \leq a + b \leq \abs{a } + \abs{ b} \\
        &\iff - \paren{\abs{a } + \abs{b}} \leq a + b \leq \abs{a} + \abs{b}
    \end{align*}

    If $\abs{a}+\abs{b}=0$, then the displayed bounds give $a+b=0$, so the
    claim is immediate. Otherwise, part (c) of the preceding theorem gives
    \begin{equation*}
        \abs{a + b } \leq \abs{a} + \abs{b}.
    \end{equation*}
\end{proof}

\begin{corollary}
    \begin{enumerate}
        \item $\abs{\abs{a } - \abs{b}} \leq \abs{a - b}. $
        \item $\abs{a - b } \leq \abs{a } + \abs{b}$
    \end{enumerate}
\end{corollary}

\begin{proof}[Proof of (1)]
    Let $a, b \in \R$. We know
    \begin{align*}
        a = a - b + b &\implies \abs{a} = \abs{(a - b) + b} \leq \abs{a - b} + \abs{b} \\
        &\implies \abs{a } - \abs{b } \leq \abs{a - b}
    \end{align*}

    \begin{align*}
        b = b - a + a &\implies \abs{b } = \abs{(b - a) + a } \leq \abs{b - a } + \abs{a }\\
        &\implies \abs{b } - \abs{a } \leq \abs{b-a}
            = \abs{-(a-b)}=\abs{a-b}. \\
        &\implies \abs{a } - \abs{b } \geq - \abs{a - b}
    \end{align*}

    Thus,
    \begin{equation*}
        -\abs{a - b } \leq \abs{a } - \abs{b } \leq \abs{a - b }
    \end{equation*}

    Therefore,
    \begin{equation*}
        \abs{\abs{a } - \abs{b }} \leq \abs{a - b}
    \end{equation*}
\end{proof}

\begin{proof}[Proof of (2)]
    Let $a, b \in \R$. We know
    \begin{align*}
        \abs{a - b} = \abs{a + \paren{-b }} \leq \abs{a} + \abs{-b }
        = \abs{a } + \abs{b}.
    \end{align*}
\end{proof}

%!TEX root=../../note
\subsection{Bounds}
Bounds make sense in any ordered set. The special feature of $\R$ is that every nonempty set bounded above has a supremum in $\R$; this need not hold in $\Q$.

\begin{definition}
    [Bounded Sets]
    Let $S$ be a nonempty subset of $\R$. The set $S$ is said to be 
    \begin{enumerate}
        \item bounded above: there exists $u \in \R$ such that $x \leq u$ for all $x \in S$. 
        \item bounded below: there exists $w \in \R$ such that $w \leq x$ for all $x \in S$. 
        \item bounded: it is both bounded above and below. 
        \item unbounded: not bounded. 
    \end{enumerate}
\end{definition}

If there exists an upper bound $u_0$ of $S$ such that whenever $u$ is an upper bound for $S$
we have $u_0 \leq u$, then $u_0$ is called the least upper bound, or the supremum, of $S$.
\begin{equation*}
    u_0 := \sup S
\end{equation*}

If there exists a lower bound $w_0$ of $S$ such that whenever $w$ is a lower bound for $S$
we have $w_0 \geq w$, then $w_0$ is called the greatest lower bound, or the infimum, of $S$.
\begin{equation*}
    w_0 := \inf S
\end{equation*}

\begin{remark}
    There can be only one supremum/infimum of a given set $S$ of $\R$.

    Idea: Suppose $u_1$ and $u_2$ are both suprema of $S$.

    Since $u_1$ is a supremum and $u_2$ is an upper bound, we get
    \begin{equation*}
        u_1 \leq u_2.
    \end{equation*}

    Since $u_2$ is a supremum and $u_1$ is an upper bound, we get
    \begin{equation*}
        u_2 \leq u_1.
    \end{equation*}

    Thus, $u_1 = u_2$.
\end{remark}


\begin{lemma}
    An upper bound $u$ of a nonempty set $S$ in $\R$ is the supremum of $S$ if and only if
    \begin{equation*}
        (\forall \epsilon > 0)(\exists x_\epsilon \in S)u - \epsilon < x_\epsilon.
    \end{equation*}
\end{lemma}

\begin{proof}
    First suppose $u=\sup S$. Let $\epsilon>0$. Since
    \begin{equation*}
        u-\epsilon<u,
    \end{equation*}
    the number $u-\epsilon$ cannot be an upper bound: otherwise it would be an upper bound smaller than $u$. By the definition of upper bound, there is therefore some $x_\epsilon\in S$ such that
    \begin{equation*}
        u-\epsilon<x_\epsilon\leq u.
    \end{equation*}

    Conversely, suppose $u$ is an upper bound and that for every $\epsilon>0$ there exists $x_\epsilon\in S$ with $u-\epsilon<x_\epsilon$. If $u$ were not the least upper bound, there would be another upper bound $w<u$. Choose
    \begin{equation*}
        \epsilon=u-w>0.
    \end{equation*}
    The assumed condition then yields
    \begin{equation*}
        w=u-\epsilon<x_\epsilon,
    \end{equation*}
    contradicting that $w$ is an upper bound. Hence $u=\sup S$.
\end{proof}

\begin{example}
    Consider the set $S = \set{x \in \R : 0 \leq x \leq 1}$.
    \begin{itemize}
        \item $\sup S = 1$.
        \begin{proof}
        Every $x\in S$ satisfies $x\leq1$, so $1$ is an upper bound. To see directly that no smaller number works, let $u<1$. If $u<0$, then $0\in S$ lies above $u$. If $0\leq u<1$, set $x=(u+1)/2$. Then
        \begin{equation*}
            u<x=\frac{u+1}{2}<1,
        \end{equation*}
        so $x\in S$ lies above $u$. Thus every number below $1$ fails to be an upper bound, and $\sup S=1$.
        \end{proof}
        \item $\inf S = 0$.
        \begin{proof}
            Since $x\geq0$ for all $x\in S$, $0$ is a lower bound. If $w>0$, then $0\in S$ satisfies $0<w$, so $w$ is not a lower bound. Thus no number greater than $0$ is a lower bound, and $\inf S=0$.
        \end{proof}
    \end{itemize}
\end{example}

%!TEX root=../../note.tex
\subsection{The Completeness Property of $\R$}
\begin{axiom}[Completeness Property]
    Every nonempty set of real numbers that has an upper bound also has a supremum in $\R$.
\end{axiom}

\begin{corollary}
    Every nonempty set of real numbers that has a lower bound also has an infimum in $\R$.
\end{corollary}
\begin{proof}
    Let $S \subseteq \R$ be nonempty with a lower bound $w$, that is, $w \leq s$ for all $s \in S$. As a consequence, $-w \geq -s$ for all $s \in S$.

    Define $S' = \set{-s : s \in S}$. Then $S'$ is nonempty, and
    \begin{equation*}
        -w \geq -s, \ \forall s \in S \implies -w \text{ is an upper bound for } S'.
    \end{equation*}

    By the Completeness Axiom, $S'$ has a supremum. Say
    \begin{equation*}
        u := \sup S' \in \R.
    \end{equation*}

    We claim that $-u = \inf S$.
    \begin{itemize}
        \item \textbf{Lower bound:} Since $u = \sup S' \geq -s$ for all $s \in S$, we get $-u \leq s$ for all $s \in S$. So $-u$ is a lower bound for $S$.
        \item \textbf{Greatest lower bound:} Let $w$ be any lower bound for $S$. As shown above, $-w$ is then an upper bound for $S'$, so $u = \sup S' \leq -w$, which gives $w \leq -u$.
    \end{itemize}
    Thus $-u$ is a lower bound for $S$ that is at least as large as every other lower bound, so $-u = \inf S$.
\end{proof}

\begin{remark}
    Recall the sup/inf recipe from the previous section: to verify $u = \sup S$, check (i) $u$ is an upper bound for $S$, and (ii) no number smaller than $u$ is an upper bound for $S$ (equivalently, the $\epsilon$-characterization: for every $\epsilon>0$ there is $x_\epsilon\in S$ with $u-\epsilon<x_\epsilon$). We will use this recipe repeatedly below, most notably in the Archimedean Property and in the existence of $\sqrt2$.
\end{remark}

\begin{example}
    \begin{equation*}
        S = \set{\paren{1 + \frac{1}{n }}^n : n \in \N}
    \end{equation*}
    
    For every $n\in\N$, the binomial theorem gives
    \begin{equation*}
        \paren{1+\frac1n}^n
        =1+\binom n1\frac1n+\sum_{k=2}^n\binom nk\frac1{n^k}>1.
    \end{equation*}
    Indeed, for $0\leq k\leq n$,
    \begin{equation*}
        \binom nk=\frac{n(n-1)\cdots(n-k+1)}{k!}\leq\frac{n^k}{k!}.
    \end{equation*}
    For $k\geq2$, we have
    $k!\geq2^{k-1}$, because each factor from $2$ through $k$ is at least $2$. Therefore
    \begin{equation*}
        \paren{1+\frac1n}^n
        \leq\sum_{k=0}^n\frac1{k!}
        < 1+1+\sum_{k=2}^{\infty}\frac1{2^{k-1}}=3.
    \end{equation*}
    Thus $1<\paren{1+1/n}^n<3$ for all $n\in\N$.
\end{example}

\subsubsection{The Archimedean Property}
We now prove that $\N$ is not bounded above in $\R$. The proof combines two ingredients:
\begin{itemize}
    \item Completeness of $\R$: this is what lets us form $u = \sup \N$ in the first place. Note carefully that $u$ here is the supremum of $\N$ \emph{regarded as a subset of $\R$} — completeness says nothing about whether $u$ itself is a natural number, and indeed the entire argument below is designed to show that assuming $\N$ is bounded above (so that such a $u$ exists) leads to a contradiction.
    \item The inductive property of $\N$ (that $\N$ is closed under $+1$): this is the purely arithmetic fact that lets us manufacture a natural number strictly larger than $u$, which is what breaks the supremum.
\end{itemize}

\begin{definition}
    [Archimedean Property]
    If $x \in \R$, then there exists $n_x \in \N$ such that $x \leq n_x$. 
\end{definition}

\begin{remark}[Proof idea]
    We argue by contradiction: suppose some real number $x$ exceeds every natural number, so that $\N$ is bounded above by $x$. Completeness then hands us $u=\sup\N\in\R$. The key move is that $u-1$, being strictly less than $u$, cannot itself be an upper bound for $\N$ (since $u$ is the \emph{least} upper bound) — this forces some natural number $m$ to lie strictly beyond $u-1$. But then $m+1$ is again a natural number (by the inductive property) sitting strictly beyond $u$, which contradicts $u$ being an upper bound for $\N$ in the first place. So no such $x$ can exist.
\end{remark}
\begin{proof}
    Suppose, for contradiction, that $\exists x \in \R$ such that $x > n$ for all $n \in \N$.

    Then $x$ is an upper bound for $\N$. Since $1 \in \N$, $\N \neq \emptyset$, so by the Completeness Axiom, $\N$ has a supremum in $\R$.

    Say $u = \sup \N$. Since
    \begin{equation*}
        u - 1 < u,
    \end{equation*}
    and $u$ is the \emph{least} upper bound of $\N$, no number strictly smaller than $u$ can be an upper bound for $\N$ (an upper bound smaller than the least upper bound would contradict its minimality). Applying this general principle to $u-1$, we conclude that $u - 1$ is not an upper bound for $\N$. Hence
    \begin{equation*}
        \exists m \in \N \text{ such that } u - 1 < m,
    \end{equation*}
    which gives
    \begin{equation*}
        u < m + 1.
    \end{equation*}
    But $m + 1 \in \N$ (since $\N$ is closed under $+1$), and $u$ is an upper bound for $\N$, so $m + 1 \leq u$. This contradicts $u < m+1$: we have simultaneously $u < m+1$ and $m+1 \leq u$.

    Therefore no such $x$ exists; that is, for every $x \in \R$ there exists $n_x \in \N$ with $x \leq n_x$.
\end{proof}

\[
    \boxed{\text{Completeness} + \text{closure of }\N\text{ under }+1 \Longrightarrow \text{Archimedean Property}}
\]

\begin{remark}
    The step ``$u-1<u$, so $u-1$ is not an upper bound'' is exactly the $\epsilon$-characterization of the supremum from the previous section, applied with $\epsilon=1$: if $u=\sup S$, then for every $\epsilon>0$, $u-\epsilon$ fails to be an upper bound for $S$, so some element of $S$ exceeds $u-\epsilon$. Here $S=\N$ and $\epsilon=1$ give the natural number $m$ with $u-1<m$.
\end{remark}

\begin{corollary}
    The following are consequences of the Archimedean Property:
    \begin{enumerate}[label=(\alph*)]
        \item If
        \[
            S = \left\{ \frac{1}{n} \mid n \in \mathbb{N} \right\},
        \]
        then
        \[
            \inf S = 0.
        \]

        \item If $t > 0$, there exists $n_t \in \mathbb{N}$ such that
        \[
            0 < \frac{1}{n_t} < t.
        \]

        \item If $y > 0$, there exists $n_y \in \mathbb{N}$ such that
        \[
            n_y - 1 \leq y < n_y.
        \]
    \end{enumerate}
\end{corollary}

\begin{remark}[Proof idea]
    Part (a) uses the sup/inf recipe: $0$ is clearly a lower bound for $S$, and part (b) shows no larger number can be, since $1/n$ can be made smaller than any prescribed $t>0$. Part (b) is a direct rescaling of the Archimedean Property: bound $1/t$ above by some natural number $m$, then $n_t=m+1$ works. Part (c) is the most subtle: we want the \emph{first} natural number strictly to the right of $y$, so we collect \emph{all} candidates into a set $A=\set{n\in\N:y<n}$ and take the least element of $A$ — this is a well-ordering argument, not a supremum argument (see the remark following the proof for why).
\end{remark}
\begin{proof}
    \begin{enumerate}[label=(\alph*)]
        \item Clearly $0$ is a lower bound for $S$. Let $b>0$. By part (b), there exists $n\in\N$ such that $1/n<b$. Since $1/n\in S$, this means $b$ is not a lower bound for $S$. As $b>0$ was arbitrary, no number greater than $0$ is a lower bound for $S$. Combined with $0$ being a lower bound, $\inf S=0$.

        \item By the Archimedean Property, choose $m\in\N$ with $1/t\leq m$. Then $n_t=m+1\in\N$ satisfies $n_t>m\geq1/t$, so $n_t>1/t$, and dividing by $n_t>0$ gives $0<1/n_t<t$.

        \item We want the least natural number lying strictly to the right of $y$. To find it, we collect every natural number with this property into the set
        \[
            A=\set{n\in\N:y<n},
        \]
        and then take the smallest element of $A$; this smallest element is automatically the \emph{first} natural number beyond $y$, which is precisely what we need to pin $y$ between two consecutive integers.

        By the Archimedean Property, $A$ is nonempty (some natural number exceeds $y$). Since $A\subseteq\N$ is nonempty, the Well-Ordering Principle guarantees that $A$ has a least element; let $n_y=\min A$. Then $n_y\in A$, so $y<n_y$.

        It remains to show $n_y-1\leq y$. Suppose instead that $n_y-1>y$. Since $y>0$ and $n_y\geq1$ (as $n_y\in\N$), if $n_y=1$ then $n_y-1=0>y$ is impossible because $y>0$; so $n_y\geq2$ and $n_y-1\in\N$. Then $n_y-1\in\N$ together with $n_y-1>y$ shows $n_y-1\in A$. But $n_y-1<n_y$, contradicting that $n_y$ is the \emph{least} element of $A$. Hence $n_y-1\leq y$, and combined with $y<n_y$ we obtain
        \[
            n_y-1\leq y<n_y.
        \]
    \end{enumerate}
\end{proof}

\begin{remark}
    Part (c) is a \emph{minimum} argument, not a supremum argument, and this distinction matters. The set $A=\set{n\in\N:y<n}$ is a subset of $\N$, and in general $A$ is \emph{unbounded above} in $\R$: if $y$ is not itself a boundary case, every sufficiently large natural number also lies in $A$. So $\sup A$ would not even be a finite real number, and asking for $\sup A$ would be the wrong tool here. What we actually need is the \emph{least} element of $A$, and the Well-Ordering Principle supplies this regardless of whether $A$ is bounded above — it only requires $A$ to be a nonempty subset of $\N$. This is in direct contrast to the supremum arguments used elsewhere in this section (for $\sup\N$ in the Archimedean Property, and for $\sup S$ in the existence of $\sqrt2$ below), where the relevant sets are subsets of $\R$ and boundedness above is essential.
\end{remark}

\begin{proposition}
    There exists a real number $x$ such that $x^2 = 2$.
\end{proposition}

\begin{remark}[Proof idea]
    We want to build $\sqrt2$ directly as a supremum. Consider
    \[
        S = \set{t \in \R : t \geq 0, t^2 < 2},
    \]
    the set of all nonnegative reals whose square undershoots $2$. Intuitively, $S$ consists of all the ``candidates'' for $\sqrt2$ from below, and its supremum $x=\sup S$ should sit exactly at the boundary between squares less than $2$ and squares greater than $2$ — that is, $x^2$ should equal $2$ precisely. To prove this rigorously, we rule out the two ways $x$ could fail to be $\sqrt2$: if $x^2<2$, then $x$ is not yet at the boundary and can be nudged slightly to the right while staying in $S$, contradicting that $x$ is an upper bound; if $x^2>2$, then $x$ has overshot the boundary and can be nudged slightly to the left while remaining an upper bound for $S$, contradicting that $x$ is the \emph{least} upper bound. Both nudges are constructed using the Archimedean Property to choose $n\in\N$ large enough. Note that $x=\sup S$ is a genuine completeness argument: $S\subseteq\R$ is bounded above, and we need the \emph{least upper bound} of a set of reals, not a well-ordering argument as in part (c) above.
\end{remark}
\begin{proof}
    Define
    \begin{equation*}
        S = \set{t \in \R : t \geq 0, t^2 < 2}.
    \end{equation*}
    We have $1 \in S$, so $S \neq \emptyset$, and $2$ is an upper bound for $S$: if $t\in S$ and $t>2$, then $t^2>4>2$, a contradiction. By the Completeness Axiom, $S$ has a supremum in $\R$.

    Let $x = \sup S$. Note $x \geq 1 > 0$. We show $x^2 = 2$ by ruling out $x^2 < 2$ and $x^2 > 2$.

    \textbf{Case 1: $x^2 < 2$.} We find $n \in \N$ large enough that $x + \frac{1}{n} \in S$, contradicting that $x$ is an upper bound for $S$.

    \textit{Step 1 (bound the perturbed square).} Since $x \geq 1$ and $n\geq1$,
    \begin{align*}
        \paren{x + \frac{1}{n}}^2 &= x^2 + \frac{2x}{n} + \frac{1}{n^2} \\
        &\leq x^2 + \frac{2x}{n} + \frac{1}{n} \\
        &= x^2 + \frac{2x+1}{n}.
    \end{align*}

    \textit{Step 2 (choose $n$ via the Archimedean Property).} Since $x^2 < 2$, the quantity $2-x^2$ is a fixed positive number, so the Archimedean Property lets us choose $n\in\N$ large enough that
    \begin{equation*}
        \frac{2x+1}{n} < 2 - x^2.
    \end{equation*}

    \textit{Step 3 (derive the contradiction).} Combining Steps 1 and 2,
    \begin{equation*}
        \paren{x + \frac{1}{n}}^2 < x^2 + (2 - x^2) = 2,
    \end{equation*}
    so $x + \frac{1}{n} \in S$. But $x + \frac{1}{n} > x = \sup S$, contradicting that $x$ is an upper bound for $S$. Hence $x^2 \geq 2$.

    \[
        \boxed{
        \begin{array}{c}
            x^2<2 \Longrightarrow x\text{ can be moved slightly to the right} \\
            \text{and remain in }S \Longrightarrow x\neq\sup S.
        \end{array}
        }
    \]

    \textbf{Case 2: $x^2 > 2$.} We find $n \in \N$ large enough that $x - \frac{1}{n}$ is an upper bound for $S$ smaller than $x$, contradicting that $x = \sup S$.

    \textit{Step 1 (bound the perturbed square).} Since $1/n^2>0$,
    \begin{equation*}
        \paren{x - \frac{1}{n}}^2 = x^2 - \frac{2x}{n} + \frac{1}{n^2} > x^2 - \frac{2x}{n}.
    \end{equation*}

    \textit{Step 2 (choose $n$ via the Archimedean Property).} Since $x^2>2$, the quantity $x^2-2$ is a fixed positive number, so the Archimedean Property lets us choose $n\in\N$ large enough that
    \begin{equation*}
        \frac{2x}{n} < x^2 - 2.
    \end{equation*}

    \textit{Step 3 (derive the contradiction).} Combining Steps 1 and 2,
    \begin{equation*}
        \paren{x - \frac{1}{n}}^2 > x^2 - (x^2 - 2) = 2.
    \end{equation*}
    So for any $t \in S$, we have $t^2 < 2 < \paren{x - \frac{1}{n}}^2$. Also $x-1/n\geq0$ because $x\geq1$ and $n\geq1$. Since $u\mapsto u^2$ is strictly increasing on $[0,\infty)$, this implies $t < x-\frac1n$. Thus $x - \frac{1}{n}$ is an upper bound for $S$ smaller than $x = \sup S$, contradicting that $x$ is the \emph{least} upper bound. Hence $x^2 \leq 2$.

    \[
        \boxed{
        \begin{array}{c}
            x^2>2 \Longrightarrow x\text{ can be moved slightly to the left} \\
            \text{and remain an upper bound for }S \Longrightarrow x\neq\sup S.
        \end{array}
        }
    \]

    Combining both cases, $x^2 = 2$.
\end{proof}

\subsubsection{Density of $\Q$}

\begin{remark}
    Consider the set
    \[
        T=\{r\in\mathbb{Q}: r^2<2\},
    \]
    which mirrors the set $S$ used above to construct $\sqrt2$, but restricted to rational numbers. As before, $T\neq\varnothing$ and $T$ is bounded above in $\Q$: for instance $0\in T$, and every $r\in T$ satisfies $r<2$.

    If $T$ is instead regarded as a subset of $\R$, completeness of $\R$ guarantees that the supremum
    \[
        y=\sup_{\R} T
    \]
    exists. Running exactly the same case-split argument as in the proof above (with $T$ in place of $S$) shows that $y^2=2$, so $y=\sqrt2$.

    But $\sqrt2\notin\Q$, so this candidate supremum $y$ is not itself a rational number. Since $y$ was the \emph{only} possible least upper bound of $T$ in $\R$, no rational number can serve as a least upper bound of $T$ in $\Q$ either — that is, $T$ has no supremum in $\Q$.

    We have therefore produced a nonempty subset $T\subset\Q$ that is bounded above in $\Q$ but has no least upper bound in $\Q$. Thus $\Q$ is \emph{not} complete: completeness genuinely depends on $\R$ containing irrational numbers like $\sqrt2$ to serve as suprema for sets like $T$.
\end{remark}

\begin{theorem}
    If $x$ and $y$ are any real numbers with $x < y$, then there exists a rational number $r \in \Q$ such that $x < r < y$. 
\end{theorem}

\begin{remark}[Proof idea]
    We first treat the case $0<x<y$. The idea is to build a rational number of the form $m/n$ landing strictly between $x$ and $y$. First we choose $n\in\N$ large enough (via the Archimedean Property) that $1/n<y-x$; this makes consecutive multiples of $1/n$ closer together than the gap $y-x$, which guarantees that at least one multiple of $1/n$ will fall inside $(x,y)$. Then we choose $m$ to be the integer just to the right of $xn$, using the integer-location result (part (c) above) applied to $xn$. A short computation shows $xn<m<yn$, and dividing by $n$ places $r=m/n$ strictly between $x$ and $y$. The general case $x\leq0$ is reduced to the positive case by shifting both $x$ and $y$ up by a large enough integer $k$, applying the positive case, and shifting back down.
\end{remark}
\begin{proof}
    First suppose that $0<x<y$. Then $y-x>0$. By the Archimedean Property, there exists $n\in\N$ such that
    \begin{equation*}
        \frac1n<y-x.
    \end{equation*}
    This choice of $n$ ensures that consecutive multiples of $1/n$ are spaced closer together than the gap between $x$ and $y$, so some multiple of $1/n$ ought to land inside $(x,y)$; we now locate it explicitly. Multiplying $1/n<y-x$ by $n>0$ shows that
    \begin{equation*}
        xn+1<yn.
    \end{equation*}

    We now find a natural number between $xn$ and $xn+1$. Since $xn>0$, part~(c) of the preceding corollary (with $y$ there replaced by $xn$) gives some $m\in\N$ with
    \begin{equation*}
        m-1\leq xn<m.
    \end{equation*}
    Equivalently,
    \begin{equation*}
        xn<m\leq xn+1.
    \end{equation*}
    This is exactly why part (c) is invoked here: $m$ is, by construction, the first natural number strictly to the right of $xn$, so $xn<m$ holds automatically. Combining this with $xn+1<yn$ obtained above, we have
    \begin{equation*}
        xn<m<yn,
    \end{equation*}
    where the left inequality is immediate from $m>xn$, and the right inequality follows since $m\leq xn+1<yn$. Dividing throughout by $n>0$, we get
    \begin{equation*}
        x<\frac{m}{n}<y.
    \end{equation*}
    Since $m\in\Z$ and $n\in\N$, $m/n\in\Q$. This proves the theorem when $0<x<y$.

    Now suppose that $x\leq0$. Choose $k\in\N$ such that $-x<k$ (possible by the Archimedean Property). Then adding $k$ to $x<y$ and to $-x<k$ gives
    \begin{equation*}
        0<x+k<y+k.
    \end{equation*}
    By the positive case just proved, there exists $q\in\Q$ such that $x+k<q<y+k$. Let $r=q-k$. Since $q\in\Q$ and $k\in\N\subseteq\Q$, we have $r\in\Q$, and subtracting $k$ from each part of $x+k<q<y+k$ gives
    \begin{equation*}
        x<r<y.
    \end{equation*}
\end{proof}

\begin{corollary}
    If $x$ and $y$ are any real numbers with $x < y$, then there exists an irrational number $z$ such that $x < z < y$.
\end{corollary}

\begin{remark}[Proof idea]
    Rescale the interval $(x,y)$ by dividing by $\sqrt2$, obtaining $(a,b)$ with $a=x/\sqrt2$ and $b=y/\sqrt2$. Apply the density of $\Q$ to find a nonzero rational $r\in(a,b)$. Then $z=\sqrt2\,r$ lands back in $(x,y)$, and $z$ must be irrational: if it were rational, $\sqrt2=z/r$ would be a ratio of rationals and hence rational, contradicting the irrationality of $\sqrt2$.
\end{remark}
\begin{proof}
    Set $a=x/\sqrt2$ and $b=y/\sqrt2$. Since $\sqrt2>0$, we have $a<b$. Choose a nonzero rational number $r$ with $a<r<b$: if $b\leq0$, choose $r\in\Q$ with $a<r<b$; if $a\geq0$, choose $r\in\Q$ with $a<r<b$; and if $a<0<b$, choose $r\in\Q$ with $0<r<b$. In each case, the density theorem provides such an $r$, and $r\neq0$.

    Let $z=\sqrt2r$. Multiplying $a<r<b$ by $\sqrt2$ gives $x<z<y$. If $z$ were rational, then $\sqrt2=z/r$ would be rational, a contradiction. Thus $z$ is irrational.
\end{proof}

\begin{corollary}
    Let $x, y \in \R$, with $x < y$. Then there exists infinitely many rationals between $x$ and $y$. 
\end{corollary}

\begin{remark}[Proof idea]
    Apply the density theorem repeatedly: once we have a rational $r_k$ strictly between $x$ and $y$, the theorem applied to the (still nonempty) interval $(r_k,y)$ produces a new rational $r_{k+1}$ strictly between $r_k$ and $y$. Iterating this produces an infinite strictly increasing sequence of distinct rationals, all lying in $(x,y)$.
\end{remark}
\begin{proof}
    Choose $r_1\in\Q$ with $x<r_1<y$. Having chosen $r_k\in\Q$ with $x<r_k<y$, apply the density theorem to $r_k<y$ to choose $r_{k+1}\in\Q$ such that
    \[
        r_k<r_{k+1}<y.
    \]
    This produces infinitely many distinct rational numbers $r_1,r_2,\ldots$ in $(x,y)$.
\end{proof}

Are there more rational or irrational numbers? To answer this question, we need to know in what sense. 
\begin{definition}
    [Cardinality]
    A function $f:A\to B$ is a \emph{bijection} if it is both injective (one-to-one) and surjective (onto). We say that $A$ and $B$ have the same \emph{cardinality}, written
    \[
        |A|=|B|,
    \]
    if there exists a bijection $f:A\to B$. Thus, cardinality records the size of a set up to a bijection.
\end{definition}

\begin{definition}
    A set $S$ is \emph{finite} if $S=\varnothing$ or if there is a bijection
    \[
        \{1,2,\ldots,n\}\to S
    \]
    for some $n\in\N$. A set $S$ is \emph{countably infinite} if there is a bijection $f:\N\to S$.
\end{definition}
\begin{example}
    \begin{enumerate}[label = (\alph*)]
        \item The function $f:\N\to\{2,4,6,\ldots\}$ defined by $f(n)=2n$ is a bijection. Thus the set of even natural numbers has the same cardinality as $\N$.
        \item The function $g:\N\to\Z$ defined by
        \[
            g(1)=0,\qquad g(2k)=k,\qquad g(2k+1)=-k \quad (k\in\N)
        \]
        is a bijection. Thus $\Z$ has the same cardinality as $\N$.
    \end{enumerate}
\end{example}
\begin{definition}
    A set $S$ is \emph{countable} if it is finite or countably infinite. A set is \emph{uncountable} if it is not countable.
\end{definition}

%!TEX root=../../../note
\subsection{Limits}
A sequence of real numbers is a function defined on the set $N = \set{1, 2, 3, \cdots}$ of natural numbers, whose range is contained 
in the set $\R$ of real numbers. 
\begin{equation*}
    X: \N \to \R
\end{equation*}
\begin{equation*}
    x_n: \text{ element of the sequence. }
\end{equation*}
Notation: $\set{x_n }, \set{x_n }_{n = 1}^{\infty}$

\begin{definition}
    A sequence $\set{x_n }$ in $\R$ is said to converge to $a \in \R$, or $a$ is said to be a limit of $\set{x_n }$ 
    , if for every $\epsilon > 0$ there exists a natural number $n_\epsilon \in \N$ such that for all $n \geq n_\epsilon$
    the term $x_n$ satisfy 
    \begin{equation*}
        \abs{x_n - a} < \epsilon. 
    \end{equation*} 

    The notation is 
    \begin{equation*}
        \lim_{n \to \infty} x_n = a. 
    \end{equation*}
\end{definition}

\begin{remark}[Reading the definition]
    In plain language: for every tolerance $\epsilon > 0$, all sufficiently late terms of the sequence lie inside the interval
    $\paren{a - \epsilon, a + \epsilon}$. Here $\epsilon$ is the tolerance, or precision level, we demand of the approximation
    $x_n \approx a$; the index $n_\epsilon$ tells us how late ``sufficiently late'' means, and in general depends on $\epsilon$.
    The phrase ``for all $n \geq n_\epsilon$'' says that once the sequence enters the $\epsilon$-band at index $n_\epsilon$, it
    stays inside that band for every later index -- it does not escape and re-enter. Since $n_\epsilon$ depends on $\epsilon$,
    tightening the tolerance (taking $\epsilon$ smaller) typically forces $n_\epsilon$ to be larger: the sequence may need to
    wait longer before it settles into a narrower band around $a$.
\end{remark}

The same ``eventually close'' idea appears in the $\epsilon$--$\delta$ definition of a function limit. Consider the oscillating function
\[
    f(x)=x\sin\paren{\frac1x}\qquad (x\ne0).
\]
Its oscillations become more frequent near $0$, but their amplitude is at most $|x|$, so the graph is squeezed toward the limiting value $0$.

\begin{center}
\begin{tikzpicture}[x=2.4cm,y=2.2cm,>=Stealth]
    % The epsilon-band and the delta-neighbourhood.
    \fill[red!8] (-0.9,-0.35) rectangle (0.9,0.35);
    \draw[densely dashed,red!75!black] (-0.9,0.35) -- (0.9,0.35) node[right] {$\epsilon$};
    \draw[densely dashed,red!75!black] (-0.9,-0.35) -- (0.9,-0.35) node[right] {$-\epsilon$};
    \draw[densely dashed,blue!70!black] (-0.35,-0.85) -- (-0.35,0.85) node[above] {$a-\delta$};
    \draw[densely dashed,blue!70!black] (0.35,-0.85) -- (0.35,0.85) node[above] {$a+\delta$};

    % Axes and the graph of x sin(1/x).
    \draw[->] (-0.95,0) -- (0.98,0) node[right] {$x$};
    \draw[->] (0,-0.9) -- (0,0.93) node[above] {$y$};
    \draw[very thick,teal!70!black,domain=0.025:0.9,samples=600,smooth]
        plot (\x,{\x*sin((1/\x) r)});
    \draw[very thick,teal!70!black,domain=-0.9:-0.025,samples=600,smooth]
        plot (\x,{\x*sin((1/\x) r)});
    \fill (0,0) circle (0.7pt);
    \node[below] at (0,-0.12) {$a=L=0$};
    \node[teal!70!black,anchor=south west] at (0.47,0.67) {$f(x)=x\sin(1/x)$};
    \node[blue!70!black,below] at (0,-0.88) {$|x-a|<\delta$};
\end{tikzpicture}
\end{center}

\begin{remark}[Proof idea]
    We want to control $|f(x)-L| = \abs{x\sin\paren{1/x}}$. Since $|\sin(\theta)| \leq 1$ for every real $\theta$, the factor
    $\sin(1/x)$ never contributes more than a factor of $1$, so $|f(x)| \leq |x|$: the size of $f(x)$ is bounded directly by
    the size of $x$ itself. This means we can make $|f(x)-L|$ small simply by making $|x-a|=|x|$ small, with no loss in the
    bound. Choosing $\delta=\epsilon$ then closes the chain of inequalities exactly, with no slack to spare.
\end{remark}

For a given $\epsilon>0$, choose $\delta=\epsilon$. If
\[
    0<|x-a|=|x|<\delta,
\]
then
\[
    |f(x)-L|=\abs{x\sin\paren{\frac1x}}
    \leq |x| \qquad \text{(since } |\sin(\theta)| \leq 1 \text{ for all } \theta \in \R \text{)}
    <\delta=\epsilon.
\]
Thus every $x$ in the $\delta$-neighbourhood of $a=0$ (apart from $a$ itself) has its function value inside the $\epsilon$-band around $L=0$.

%!TEX root=../note.tex
\section{Sequences and Limits}
%!TEX root=../../note.tex

\subsection{Convergent Sequences}
\subsubsection{The Limit of a Sequence}

A sequence of real numbers is a function $X:\N\to\R$. We write $x_n=X(n)$ for its $n$th term and denote the sequence by $\set{x_n}_{n=1}^{\infty}$.

\begin{definition}
    A sequence $\set{x_n}$ in $\R$ converges to $a\in\R$ if, for every $\epsilon>0$, there is an index $n_\epsilon\in\N$ such that
    \[
        n\geq n_\epsilon \quad\Longrightarrow\quad \abs{x_n-a}<\epsilon.
    \]
    We then write $\lim_{n\to\infty}x_n=a$.
\end{definition}

In symbols,
\[
    x_n\to a \iff (\forall\epsilon>0)(\exists n_\epsilon\in\N)(\forall n\geq n_\epsilon)\,\bigl[\abs{x_n-a}<\epsilon\bigr].
\]
In words: for every tolerance $\epsilon>0$, all sufficiently late terms lie in $(a-\epsilon,a+\epsilon)$. The index $n_\epsilon$ may depend on $\epsilon$, and once $n\geq n_\epsilon$ the inequality must hold for \emph{every} later term, even if earlier terms lie outside the interval.

The points below show $x_n=1/n$. For the illustrated tolerance $\epsilon=0.3$, every term from $n_\epsilon=4$ onward lies in the shaded band around $a=0$.

\begin{center}
\begin{tikzpicture}[x=0.55cm,y=2.4cm,>=Stealth]
    \fill[red!8] (0,0) rectangle (11,0.3);
    \draw[densely dashed,red!75!black] (0,0.3) -- (11,0.3) node[right] {$a+\epsilon$};
    \draw[->] (0,0) -- (11.5,0) node[right] {$n$};
    \draw[->] (0,0) -- (0,1.15) node[above] {$x_n$};
    \node[left] at (0,0) {$a=0$};
    \draw[densely dashed,blue!70!black] (4,0) -- (4,0.95);
    \node[below] at (4,0) {$n_\epsilon=4$};
    \foreach \n in {1,...,10} {
        \fill[teal!70!black] (\n,{1/\n}) circle (1.7pt);
    }
    \node[teal!70!black,anchor=west] at (1.3,0.85) {$x_n=1/n$};
\end{tikzpicture}
\end{center}

For $a\in\R$, the $\epsilon$-neighbourhood of $a$ is
\[
    V_\epsilon(a)=\set{x\in\R:\abs{x-a}<\epsilon}=(a-\epsilon,a+\epsilon).
\]
Thus $x_n\to a$ precisely when, for every $\epsilon>0$, there is $n_\epsilon\in\N$ such that $x_n\in V_\epsilon(a)$ for all $n\geq n_\epsilon$.

\begin{example}
    Prove that $\lim_{n\to\infty}\frac1n=0$.
\end{example}
\mdfsetup{nobreak=true}
\begin{proof}
    Let $\epsilon>0$. By the Archimedean Property, choose $n_\epsilon\in\N$ with $1/n_\epsilon<\epsilon$. For every $n\geq n_\epsilon$,
    \[
        \abs{\frac1n-0}=\frac1n\leq\frac1{n_\epsilon}<\epsilon.
    \]
    Hence $1/n\to0$.
\end{proof}
\mdfsetup{nobreak=false}

\begin{example}
    Prove that $\lim_{n\to\infty}\frac{4n+3}{n+1}=4$.
\end{example}
\begin{scratch}
    Simplify the distance first, then see how large $n$ must be:
    \[
        \abs{\frac{4n+3}{n+1}-4}=\abs{\frac{-1}{n+1}}=\frac1{n+1}<\frac1n.
    \]
    So it is enough that $\frac1n<\epsilon$. By Archimedes pick $n_\epsilon$ with $\frac1{n_\epsilon}<\epsilon$; every $n\geq n_\epsilon$ then works.
\end{scratch}
\begin{proof}
    Let $\epsilon>0$. By the Archimedean Property, choose $n_\epsilon\in\N$ with $1/n_\epsilon<\epsilon$. For every $n\geq n_\epsilon$,
    \[
        \abs{\frac{4n+3}{n+1}-4}
        =\frac1{n+1}\leq\frac1n\leq\frac1{n_\epsilon}<\epsilon.
    \]
    Therefore the sequence converges to $4$.
\end{proof}

\subsubsection{Uniqueness of the Limit}
\begin{proposition}
    If $x,y\in\R$ and $\abs{x-y}<\epsilon$ for every $\epsilon>0$, then $x=y$.
\end{proposition}
\begin{proof}
    Suppose $x\ne y$. Then $\epsilon=\abs{x-y}/2>0$, but $\abs{x-y}>\epsilon$, contradicting the hypothesis.
\end{proof}

\begin{theorem}
    A sequence in $\R$ can have at most one limit.
\end{theorem}
\begin{scratch}
    By the preceding proposition, it suffices to show $\abs{a-b}<\epsilon$ for every $\epsilon>0$. The only link between $a$ and $b$ is the sequence, so pass through a term $x_n$:
    \[
        \abs{a-b}\leq\abs{a-x_n}+\abs{x_n-b}.
    \]
    To make the sum $<\epsilon$, make each piece $<\epsilon/2$: that needs $n\geq n_1$ for the first and $n\geq n_2$ for the second, so take $n\geq\max\set{n_1,n_2}$.
\end{scratch}
\begin{proof}
    Suppose $x_n\to a$ and $x_n\to b$. Given $\epsilon>0$, choose $n_1,n_2\in\N$ so that $\abs{x_n-a}<\epsilon/2$ whenever $n\geq n_1$ and $\abs{x_n-b}<\epsilon/2$ whenever $n\geq n_2$. For $n\geq\max\set{n_1,n_2}$, the triangle inequality gives
    \[
        \abs{a-b}\leq\abs{a-x_n}+\abs{x_n-b}<\frac\epsilon2+\frac\epsilon2=\epsilon.
    \]
    This holds for every $\epsilon>0$, so the preceding proposition gives $a=b$.
\end{proof}

\subsubsection{Divergence and Boundedness}
A sequence diverges if it does not converge to any real number. For example, consider $\set{x_n}=\set{0,1,0,1,\ldots}$, with $x_n=0$ for odd $n$ and $x_n=1$ for even $n$.

Both $0$ and $1$ keep appearing arbitrarily late, so the terms cannot all remain close to one proposed limit.

Suppose $x_n\to a$ and take $\epsilon=1/2$. There is an $n_\epsilon$ such that $\abs{x_n-a}<1/2$ for every $n\geq n_\epsilon$. Choose an odd and an even index beyond $n_\epsilon$. The odd term gives $\abs{a}<1/2$, hence $-1/2<a<1/2$. The even term gives $\abs{1-a}<1/2$, hence $1/2<a<3/2$. These intervals do not overlap, a contradiction. Thus this sequence diverges.

\begin{definition}
    A sequence $\set{x_n}$ is bounded if there is an $M\in\R$ such that $\abs{x_n}\leq M$ for every $n\in\N$.
\end{definition}

\begin{theorem}
    Every convergent sequence of real numbers is bounded.
\end{theorem}
\begin{idea}
    The tail stays within $1$ of the limit; only finitely many terms come before it, and a finite set has a maximum.
\end{idea}
\begin{proof}
    Suppose $x_n\to a$. Taking $\epsilon=1$, choose $n_1\in\N$ so that $\abs{x_n-a}<1$ for all $n\geq n_1$. For these terms, the triangle inequality gives
    \[
        \abs{x_n}\leq\abs{x_n-a}+\abs{a}<1+\abs{a}.
    \]
    The remaining terms form a finite set. Thus
    \[
        M=\max\set{\abs{x_1},\ldots,\abs{x_{n_1-1}},1+\abs{a}}
    \]
    bounds every term of the sequence (when $n_1=1$, take $M=1+\abs{a}$).
\end{proof}

The converse fails: $\set{0,1,0,1,\ldots}$ is bounded but, as shown above, diverges.


\begin{example}
    Consider the following sequences. 
    \begin{itemize}
        \item $\set{a_n } = \set{\frac{n^2 + 1 }{n }}$
        
        Take a look at this sequence: 
        \begin{equation*}
            \abs{\frac{n^2 + 1 }{n }} = \abs{n + \frac{1}{n }} = n + \frac{1}{n}
        \end{equation*}
    
        For any $M > 0$, there exist $n_m \in \N$ such that $n_M > M$.  Thus, 
        \begin{equation*}
            \abs{a_{n_M}} > M. 
        \end{equation*}

        \item $\set{a_n } = \set{\paren{-1}^n n}$
        
        Take a look at this sequence:

        \begin{equation*}
            \abs{a_n } = \abs{\paren{-1}^n n } = \abs{\paren{-1}^n }\abs{n } = \abs{n } > M
        \end{equation*}
    \end{itemize}
\end{example}

\subsubsection{Divergence to $\pm\infty$}
\begin{definition}
    A sequence $\set{x_n }$ diverges to $+\infty$, that is, $\lim_{x_n } = \infty$, if for all real number $M > 0$, there exists $n_M \in \N$
    such that $a_n > M$ for all $n \geq n_M$. 
\end{definition}

\begin{definition}
    A sequence $\set{x_n }$ diverges to $-\infty$, that is, $\lim_{x_n } = -\infty$, if for all real number $M > 0$, there exists $n_M \in \N$
    such that $a_n < - M$ for all $n \geq n_M$. 
\end{definition}
We write $\lim_{n\to\infty}x_n=\pm\infty$, but such a sequence is still divergent: it does not converge to any real number.


\begin{example}
    Prove that $\lim_{n \to \infty } \frac{n^2}{1 - n} = -\infty$. 
\end{example}
\begin{proof}
    Let $M > 0$. 
\end{proof}

\begin{example}
    Prove that $\lim_{x \to \infty} \sqrt{n } + 1 = + \infty$. 
\end{example}
\begin{proof}
    The goal is to let $\sqrt{n } > m - 1$
\end{proof}

\begin{theorem}
    If $\lim_{n \to \infty} x_n = a$ and $\lim_{n \to \infty} y_n = b$, then 
    \begin{enumerate}
        \item $(\forall c \in \R)\,[\lim_{n \to \infty} c = c]$
        \item $(\forall c \in \R)\,[\lim_{n \to \infty} cx_n = ca]$
    \end{enumerate}
\end{theorem}
%!TEX root=../../note.tex
% Source: MATH 147 lecture notes; limits and order, squeeze theorem, monotone convergence.
\subsection{Limits, Order, and Monotone Sequences}
Limits preserve order, and completeness turns a bounded monotone sequence into a convergent one. These facts let us find limits even when the limit laws alone do not apply.

\subsubsection{Limits and Inequalities}
\begin{theorem}\label{thm:limit-nonneg}
    If $\lim_{n\to\infty}x_n=a$ and $x_n\geq0$ for every $n\in\N$, then $a\geq0$.\index{limit!preserves order}
    \insymbols{\bigl[\lim_{n\to\infty}x_n=a\wedge(\forall n\in\N)\,[x_n\geq0]\bigr]\Rightarrow a\geq0}
\end{theorem}
\begin{idea}
    If $a<0$, convergence would force every sufficiently late term into a small interval around $a$ that lies below zero.
\end{idea}
\begin{proof}
    Suppose instead that $a<0$. Take $\epsilon=-a/2>0$. Since $\lim_{n\to\infty}x_n=a$, some $n_\epsilon\in\N$ satisfies $\abs{x_n-a}<\epsilon$ whenever $n\geq n_\epsilon$. For such $n$,
    \[
        x_n<a+\epsilon=\frac a2<0,
    \]
    contradicting $x_n\geq0$. Thus $a\geq0$.
\end{proof}

The same conclusion holds if $x_n\geq0$ only for all sufficiently large $n$: the proof uses just one term after both thresholds.

\begin{theorem}\label{thm:limits-and-order}
    If $\lim_{n\to\infty}x_n=a$, $\lim_{n\to\infty}y_n=b$, and $x_n\leq y_n$ for every $n\in\N$, then $a\leq b$.
    \insymbols{\bigl[\lim_{n\to\infty}x_n=a\wedge \lim_{n\to\infty}y_n=b\wedge(\forall n\in\N)\,[x_n\leq y_n]\bigr]\Rightarrow a\leq b}
\end{theorem}
\begin{proof}
    Set $z_n=y_n-x_n\geq0$. The difference law (\cref{thm:limit-laws}, part (iv)) gives $\lim_{n\to\infty}z_n=b-a$. By \cref{thm:limit-nonneg}, $b-a\geq0$, so $a\leq b$.
\end{proof}

Strict inequalities are not preserved: $\frac1n>0$ for every $n\in\N$, yet $\lim_{n\to\infty}\frac1n=0$. From $x_n<y_n$ for all $n$ we may conclude only $a\leq b$.

\subsubsection{The Squeeze Theorem}
\begin{theorem}[Squeeze Theorem]\label{thm:squeeze}
    Suppose $x_n\leq y_n\leq z_n$ for every $n\in\N$ and $\lim_{n\to\infty}x_n=a$, $\lim_{n\to\infty}z_n=a$. Then $\lim_{n\to\infty}y_n=a$.\index{squeeze theorem}
    \insymbols{\bigl[(\forall n\in\N)\,[x_n\leq y_n\leq z_n]\wedge \lim_{n\to\infty}x_n=a\wedge \lim_{n\to\infty}z_n=a\bigr]\Rightarrow \lim_{n\to\infty}y_n=a}
\end{theorem}
\begin{proof}
    Let $\epsilon>0$. Choose $n_1,n_2\in\N$ such that $\abs{x_n-a}<\epsilon$ for $n\geq n_1$ and $\abs{z_n-a}<\epsilon$ for $n\geq n_2$. If $n\geq\max\set{n_1,n_2}$, then
    \[
        a-\epsilon<x_n\leq y_n\leq z_n<a+\epsilon.
    \]
    Thus $\abs{y_n-a}<\epsilon$ for every such $n$, so $\lim_{n\to\infty}y_n=a$.
\end{proof}

\begin{example}\label{ex:sin-over-sqrt}
    Prove that $\lim_{n\to\infty}\frac{\sin n}{\sqrt n}=0$.
\end{example}
\begin{proof}
    Since $-1\leq\sin n\leq1$ and $\sqrt n>0$,
    \[
        -\frac1{\sqrt n}\leq\frac{\sin n}{\sqrt n}\leq\frac1{\sqrt n}.
    \]
    The outer sequences tend to $0$: given $\epsilon>0$, the Archimedean Property (\cref{def:archimedean}) supplies $n_\epsilon\in\N$ with $n_\epsilon>1/\epsilon^2$, and every $n\geq n_\epsilon$ gives $1/\sqrt n\leq1/\sqrt{n_\epsilon}<\epsilon$; the lower bound $-1/\sqrt n$ tends to $0$ by the same $n_\epsilon$. \Cref{thm:squeeze} now gives the result.
\end{proof}

\subsubsection{Monotone Sequences}
\begin{definition}\label{def:monotone}
    A sequence $\set{x_n}$ is \emph{increasing} if $x_n\leq x_{n+1}$ for every $n\in\N$, and \emph{decreasing} if $x_n\geq x_{n+1}$ for every $n\in\N$. It is \emph{monotone} if it is increasing or decreasing.\index{monotone sequence}
\end{definition}

% Keep the statement and its symbolic form in one frame.
\mdfsetup{nobreak=true}
\begin{theorem}[Monotone Convergence Theorem]\label{thm:monotone-convergence}
    A monotone sequence of real numbers converges if and only if it is bounded. More precisely:\index{monotone convergence theorem}
    \begin{enumerate}
        \item[(i)] If $\set{x_n}$ is bounded and increasing, then
            \[
                \lim_{n\to\infty}x_n=\sup\set{x_n:n\in\N}.
            \]
        \item[(ii)] If $\set{x_n}$ is bounded and decreasing, then
            \[
                \lim_{n\to\infty}x_n=\inf\set{x_n:n\in\N}.
            \]
    \end{enumerate}
    \insymbols{\begin{aligned}
        &[\set{x_n}\text{ monotone}]\\
        &\quad\Rightarrow
            \bigl[(\exists a\in\R)\,[\lim_{n\to\infty}x_n=a]\iff\set{x_n}\text{ bounded}\bigr],\\
        \text{(i)}\quad &[\set{x_n}\text{ bounded and increasing}]\\
        &\quad\Rightarrow\lim_{n\to\infty}x_n=\sup\set{x_n:n\in\N},\\
        \text{(ii)}\quad &[\set{x_n}\text{ bounded and decreasing}]\\
        &\quad\Rightarrow\lim_{n\to\infty}x_n=\inf\set{x_n:n\in\N}.
    \end{aligned}}
\end{theorem}
\mdfsetup{nobreak=false}
\begin{idea}
    Completeness gives a supremum or infimum. A term lies within $\epsilon$ of that bound, and monotonicity keeps every later term there.
\end{idea}
\begin{proof}
    Every convergent sequence is bounded (\cref{thm:convergent-bounded}), proving the forward implication. For the converse, (i) and (ii) cover the two kinds of monotone sequence.

    For (i), suppose $\set{x_n}$ is bounded and increasing. The set $S=\set{x_n:n\in\N}$ is nonempty and bounded above, so completeness gives $s=\sup S$. Let $\epsilon>0$. The $\epsilon$-characterization of the supremum (\cref{lem:eps-characterization}) gives an index $n_\epsilon$ with $s-\epsilon<x_{n_\epsilon}$. For every $n\geq n_\epsilon$, monotonicity and the upper-bound property give
    \[
        s-\epsilon<x_{n_\epsilon}\leq x_n\leq s<s+\epsilon.
    \]
    Hence $\lim_{n\to\infty}x_n=s$.

    For (ii), suppose $\set{x_n}$ is bounded and decreasing. The set $S$ is nonempty and bounded below, so $t=\inf S$ exists (\cref{cor:inf-existence}). Let $\epsilon>0$. Since $t+\epsilon>t$ and $t$ is the \emph{greatest} lower bound, $t+\epsilon$ is not a lower bound of $S$; so there is an index $n_\epsilon$ with $x_{n_\epsilon}<t+\epsilon$. For $n\geq n_\epsilon$, monotonicity and the lower-bound property give
    \[
        t-\epsilon<t\leq x_n\leq x_{n_\epsilon}<t+\epsilon.
    \]
    Thus $\lim_{n\to\infty}x_n=t$.
\end{proof}

\subsubsection{A Recursively Defined Sequence}
The Monotone Convergence Theorem is most useful for sequences defined by recursion, where no formula for $x_n$ is available. It shows that a limit exists; the recursion then determines its value.

\begin{example}\label{ex:recursive-mct}
    Let $x_1=1$ and $x_{n+1}=\frac{2x_n+3}{4}$ for all $n\in\N$. Show that $\set{x_n}$ converges and find its limit.
\end{example}
The first terms are
\[
    x_1=1,\qquad x_2=\frac{2\cdot1+3}{4}=\frac54=1.25,\qquad x_3=\frac{2\cdot\frac54+3}{4}=\frac{11}8=1.375,
\]
which suggests that the sequence increases and stays below $2$.
\begin{idea}
    Prove by induction that the sequence is bounded above and increasing; the Monotone Convergence Theorem then gives a limit $a$, and letting $n\to\infty$ in the recursion gives an equation for $a$.
\end{idea}
\begin{proof}
    \textbf{Step 1: $x_n<2$ for all $n\in\N$.} We use induction on $n$. For $n=1$, $x_1=1<2$. Let $k\in\N$ and assume $x_k<2$. Then
    \[
        2x_k<4\implies 2x_k+3<7\implies x_{k+1}=\frac{2x_k+3}{4}<\frac74<2.
    \]
    By induction, $x_n<2$ for every $n\in\N$.

    \textbf{Step 2: $x_n\leq x_{n+1}$ for all $n\in\N$.} Again by induction. For $n=1$, $x_1=1<\frac54=x_2$. Let $k\in\N$ and assume $x_k\leq x_{k+1}$. Then
    \[
        2x_k+3\leq 2x_{k+1}+3\implies x_{k+1}=\frac{2x_k+3}{4}\leq\frac{2x_{k+1}+3}{4}=x_{k+2}.
    \]
    By induction, $\set{x_n}$ is increasing.

    \textbf{Step 3: the limit.} By Step 2, $x_n\geq x_1=1$ for every $n$, and by Step 1, $x_n<2$; so $\set{x_n}$ is bounded and increasing. By \cref{thm:monotone-convergence} there is $a\in\R$ with $\lim_{n\to\infty}x_n=a$.

    The shifted sequence also satisfies $\lim_{n\to\infty}x_{n+1}=a$: if $\abs{x_n-a}<\epsilon$ for all $n\geq n_\epsilon$, then also $\abs{x_{n+1}-a}<\epsilon$ for all $n\geq n_\epsilon$, since $n+1\geq n_\epsilon$. On the other hand, $x_{n+1}=\frac12x_n+\frac34$, so the limit laws (\cref{thm:limit-laws}, parts (i)--(iii)) give
    \[
        \lim_{n\to\infty}x_{n+1}=\frac12a+\frac34=\frac{2a+3}{4}.
    \]
    A limit is unique (\cref{thm:limit-unique}), so
    \[
        a=\frac{2a+3}{4}\implies 4a=2a+3\implies a=\frac32.
    \]
    Hence $\lim_{n\to\infty}x_n=\frac32$.
\end{proof}

Existence must come first. The equation $a=\frac{2a+3}{4}$ only says what the limit would be \emph{if} there is one: for $x_1=1$ and $x_{n+1}=2x_n$, the same manipulation gives $a=2a$, so $a=0$, yet $x_n=2^{n-1}$ diverges to $+\infty$. The bound $2$ in Step 1 is not special; any upper bound at least $\frac32$ (the solution of the limit equation) would work.

\subsubsection*{In practice}
\begin{itemize}
    \item \textbf{Use the squeeze theorem.} Bound a difficult sequence above and below by sequences with the same known limit; verify that both bounds hold for all sufficiently large indices.
    \item \textbf{Use monotonicity.} Prove the direction of monotonicity and find a bound. Completeness then gives the limit as the supremum or infimum of the terms.
    \item \textbf{Recursive sequences.} Compute a few terms to guess the direction and a bound; prove both by induction; apply the Monotone Convergence Theorem; then pass to the limit on both sides of the recursion and solve for the limit.
\end{itemize}


\printindex

\end{document}
